\chapter{Procesamiento de Señales}
\label{chap:signal-processing}
{\huge$\bigstar$}{\huge$\bigstar$}{\huge$\bigstar$}{\huge$\bigstar$} Experto

\section{FFT}
\begin{lispcode}
(py:run "import numpy as np; import matplotlib.pyplot as plt
fs = 1000; t = np.linspace(0, 1, fs)
senal = np.sin(2*np.pi*50*t) + np.sin(2*np.pi*120*t)
fft = np.fft.fft(senal); freqs = np.fft.fftfreq(len(senal), 1/fs)
plt.plot(freqs[:len(freqs)//2], np.abs(fft)[:len(fft)//2]); plt.savefig('fft.png')")
\end{lispcode}

\section{Filtrado}
\begin{lispcode}
(py:run "from scipy import signal as sig; import numpy as np
b, a = sig.butter(4, 0.1, 'low'); t = np.linspace(0, 1, 1000)
clean = np.sin(2*np.pi*10*t); noisy = clean + 0.5*np.random.randn(len(t))
filtered = sig.filtfilt(b, a, noisy)")
\end{lispcode}

\section{Interpolación}
\begin{lispcode}
(py:run "from scipy.interpolate import interp1d; import numpy as np
xs = np.array([0,1,2,3,4,5]); ys = np.array([0,0.8,0.9,0.1,-0.8,-1.0])
f = interp1d(xs, ys, kind='cubic'); xs_dense = np.linspace(0,5,100)")
\end{lispcode}

